% Abstract and index terms (#81, part "Abstract and introduction"; owner: 11yashiagrawal).
% Cross-validation numbers: results/metrics/fd001_rul_cap_choice.json (#55),
% fd001_ml_capped_cv.json (#61), fd001_cv_{lstm,tcn,transformer}{,_uncapped}.json,
% fd001_hybrid_value_test.json (#68). Official-test numbers arrive with #71:
% every \pending{} must be replaced before the final build.

\providecommand{\pending}[1]{\textbf{[#1]}\typeout{Placeholder: pending #1}}

\begin{abstract}
Remaining useful life (RUL) prediction estimates how many cycles an engine has
left before it fails. Phase~1 of this project built feature-based
machine-learning baselines on the NASA C-MAPSS subset FD001. Phase~2 asks
whether deep models that read the raw sensor sequence do better under the same
leakage-controlled protocol. We first choose the training cap from the
training engines alone: a health index, a hinge fit per engine and the 90th
percentile of the onset RUL rule out caps of 150 and no cap, and support 125.
We then compare an LSTM, a temporal convolutional network and a Transformer,
which share one encoder--head contract, with Random Forest, XGBoost and
Elastic Net retuned on the same capped target and selection metric, and with
two hybrids that fuse a sequence encoder with engineered features. In
engine-grouped cross-validation, the training target matters far more than the
architecture: removing the cap costs every deep encoder 13 to 15 cycles of
RMSE, while the three encoders lie within 1.7 cycles of each other (15.6 to
17.2). The LSTM and Transformer beat machine learning on raw sensors by about
three cycles, but XGBoost on the Phase~1 engineered features remains best (14.5).
Neither hybrid beats its best single branch on accuracy under a test fixed in
advance; gated fusion only lowers the variance across seeds. On the 100
official test engines, scored once after the model freeze,
\pending{\#71: final test RMSE and NASA score of the frozen models}.
\end{abstract}

\begin{IEEEkeywords}
remaining useful life, predictive maintenance, C-MAPSS, LSTM, temporal
convolutional network, Transformer, hybrid models, evaluation protocol
\end{IEEEkeywords}
